\section{Equal-prior values with a two-dimensional retained receiver}
\label{sec:equalinitial95}
The prescribed-spectrum formula in Section~\ref{sec:spectralvalue94}
uses the biased prior of the first basis experiment.  Equal priors change
not only the baseline but also the optimal fresh acquisition.  We now
solve the equal-prior problem at every pure initial acquisition when the
old receiver has dimension at most two.  The input dimension is arbitrary,
and allowing a larger fresh reference does not change the answer.
The unrestricted qubit instance consequently admits a complete
initial-spectrum formula, including an optimal protocol without a
receiver message.

Use a weighted exact second-moment family as in
\eqref{eq:basismoments90}.  With probability $1/2$ the device is
$C_y=I/d$; otherwise draw $b$ with law $w_b$ \emph{once}, and use
$E_y^b(t)=(1-t)I/d+tP_y^b$ at both calls.  Here $d\ge2$ and
$0\le t\le1$.  The initial acquisition is the prescribed pure vector
$|C_0\rangle$, with Gram matrix $\rho=C_0^*C_0\in\mathcal D_d$;
its input marginal is $\rho^{\mathsf T}$.  It is not replaced by a
classically resolved initial mixture.  All instruments, discarded calls,
public records, null histories and norm limits follow
Definition~\ref{def:rankclass92}.  Thus the complete retained old system,
not merely one of its subsystems, has dimension at most two.

For $0\le v\le1$ define
\begin{equation}\label{eq:psidef95}
 \begin{split}
 c_d(v)&=d(d-2)+(2d-1)v,\\
 \psi_d(v)&=d-1+
 \frac{(2d-1)v}{2\left(d-1+
               d\sqrt{c_d(v)/(d^2-1)}\right)}.
 \end{split}
\end{equation}
The denominator is positive, including at $d=2,v=0$.
Set
\[
 q=\max\{\lambda_{\max}(\rho),1/2\},\qquad v_\rho=4q(1-q).
\]
The pair $(q,1-q)$ is the nonzero part, with a possible zero entry,
of the least rank-two majorant \eqref{eq:rankmajorant94}.

\begin{theorem}[Equal-prior prescribed-spectrum value]
\label{thm:equalinitial95}
For the once-selected experiment and the fixed pure first acquisition
above, and every $2\le\ell\le d$, the full class
$\mathfrak R_{2,\ell}$ has exact optimum
\begin{equation}\label{eq:equalinitial95}
 P_{2,\ell}^{\rm eq}(\rho)
   =\frac12+\frac{t^2}{2d(d+1)}\psi_d(v_\rho).
\end{equation}
A recorded instrument with at most $d$ outcomes attains this value on
the prescribed initial receiver.  The same instrument may be used after
every first device label; both retained quantum registers may have
dimension two.  The fresh preparation depends only on its recorded
receiver outcome, not on the first device label.  Its Schmidt weights
are specified in Lemma~\ref{lem:ranktwoleaf95} below.

For $t>0$ the value equals the optimized rank-two benchmark precisely
when $\lambda_{\max}(\rho)\le1/2$, and decreases strictly as that
largest eigenvalue increases above $1/2$.  At $t=0$ every value is
$1/2$.  If either retained register has dimension at most one, the
exact value at every fixed $\rho$ is instead
\begin{equation}\label{eq:rankoneinitial95}
 \frac12+\frac{t^2(d-1)}{2d(d+1)}.
\end{equation}
\end{theorem}
The first assertion covers all fresh dimensions at least two while
fixing the old dimension at two.  It makes no assertion that an old
receiver of larger dimension has the same value.  The general-rank
biased formula and the equal-prior globally optimized dimension
hierarchy retain their respective quantifiers.  If $b$ is redrawn
independently at the two calls, the averaged alternative is
$\mathcal M_C\otimes\mathcal M_C$, and all equal-prior values in
this section are $1/2$ instead.

\subsection{A concavity principle for filtering}
For a Hermitian $X$ and a positive matrix $M$ on a finite-dimensional
space, let $\mathcal F_X(M)=\|\sqrt M X\sqrt M\|_1$.
\begin{lemma}[Concavity in the filter Gram]\label{lem:filterconcavity95}
For every $M\succeq0$,
\begin{equation}\label{eq:filterconcavity95}
 \mathcal F_X(M)
  =\max_{-M\preceq Z\preceq M}\tr(XZ).
\end{equation}
In particular $\mathcal F_X$ is continuous, positively homogeneous
and concave on the positive cone.  Consequently, for fixed $A\succeq0$,
\[
 B\longmapsto
 \|(\sqrt A\otimes\sqrt B)(I-dS)(\sqrt A\otimes\sqrt B)\|_1
\]
is concave.  If $A$ is diagonal, replacing $B$ by its diagonal
pinching cannot decrease this last quantity.
\end{lemma}
\begin{proof}
Every Hermitian contraction $F$ gives
$Z=\sqrt M F\sqrt M$ in the interval $[-M,M]$.
Conversely any matrix in this interval vanishes on $\ker M$.
Indeed $M\pm Z\succeq0$ have zero quadratic form on that kernel,
so both annihilate it.  On $\operatorname{supp}M$,
$F=M^{-1/2}ZM^{-1/2}$ is a Hermitian contraction; extend it by zero.
The variational formula for the trace norm and cyclicity of the trace
prove \eqref{eq:filterconcavity95}, without invertibility of $M$.
Convex combinations of feasible optimizers are feasible for the same
combination of their bounds $M$.  This proves concavity.  Homogeneity
is immediate and continuity follows from the square-root formula.

Apply the result to $M=A\otimes B$.  If $U$ is a diagonal unitary,
$U$ commutes with $A$ and $(U\otimes U)S=S(U\otimes U)$.
Thus simultaneous conjugation shows that the value at $UBU^*$ equals
that at $B$.  Averaging diagonal sign unitaries eliminates every
off-diagonal entry, and concavity proves the pinching assertion.
Pinching may increase rank; it is used below only in an
\emph{unrestricted} upper bound.  A rank-two optimizer is subsequently
constructed, rather than inferred from a rank-preserving pinching.
\end{proof}

\subsection{The exact rank-two leaf}
\begin{lemma}[A centered swap with a rank-two old Gram]
\label{lem:ranktwoleaf95}
Let $A\succeq0$ have rank at most two and trace $a>0$.  Let
$x,1-x$ be the eigenvalues of $A/a$ on a two-dimensional support,
padded if necessary, with $1/2\le x\le1$.  For every $\ell\ge2$,
\begin{equation}\label{eq:ranktwoleaf95}
 \max_{B\in\mathcal D_d^{\le\ell}}
 \|(\sqrt A\otimes\sqrt B)(I-dS)(\sqrt A\otimes\sqrt B)\|_1
                 =a\,\psi_d(4x(1-x)).
\end{equation}
The same maximum holds without a rank restriction on $B$.  It has
an optimizer commuting with $A$ and supported on its leading two
coordinates.  For $v=4x(1-x)>0$, put
\begin{equation}\label{eq:freshweights95}
 z=2x-1,\quad e=d-1,\quad
 T=2\psi_d(v)-e,\quad
 w_*=\frac{z(eT-1)}{c_d(v)}.
\end{equation}
The two fresh eigenvalues are $(1+w_*)/2,(1-w_*)/2$ in that
ordered eigenbasis.  Here $0\le w_*\le z<1$.  At rank one, an
aligned pure fresh Gram is optimal and the value is $a(d-1)$.
At $A=0$ the value is zero.
\end{lemma}
\begin{proof}
Divide by $a$ and diagonalize $A$.
Lemma~\ref{lem:filterconcavity95} permits a diagonal $B$ for the
unrestricted maximum.  If $A$ has rank one, the norm is
$(d-1)b_1+\sum_{j>1}b_j\le d-1$, attained by $b_1=1$.
For rank two, write $m=b_1+b_2$.  The part with the fresh coordinate
outside the support of $A$ is positive and has trace $1-m$;
the swap term vanishes there.  The remaining part has norm $m$
times its normalized two-coordinate value, when $m>0$.
The two-coordinate maximum is at least $d-1$: for example $B=A$
gives
\[
 (d-1)\bigl(x^2+(1-x)^2\bigr)+2d x(1-x)
                   =d-1+2x(1-x).
\]
It follows that all fresh mass may be placed in those two coordinates.
This also constructs a candidate of rank at most two and justifies
applying the unrestricted upper at every $\ell\ge2$.

Write the two fresh weights as $(1+w)/2,(1-w)/2$, $-1\le w\le1$.
The vectors $|11\rangle,|22\rangle$ give the negative diagonal
entries $(1-d)x(1+w)/2$ and $(1-d)(1-x)(1-w)/2$.
On $|12\rangle,|21\rangle$ the block is
\[
 \begin{pmatrix}
 x(1-w)/2&-\tfrac d2\sqrt{x(1-x)(1-w^2)}\\
 -\tfrac d2\sqrt{x(1-x)(1-w^2)}&(1-x)(1+w)/2
 \end{pmatrix}.
\]
Its determinant is nonpositive.  Hence the complete trace norm is
\begin{equation}\label{eq:scalarleaf95}
 F_z(w)=\frac12\left[e(1+zw)+
       \sqrt{(z-w)^2+d^2v(1-w^2)}\right],\qquad v=1-z^2.
\end{equation}
This formula includes $w=\pm1$ by continuity.

Here is an algebraic upper certificate, which also evaluates the
maximum.  Set $C=c_d(v)>0$ and
$D(w)=1+(d^2-1)v-2zw+(1-d^2v)w^2$.
Substitution of \eqref{eq:psidef95}--\eqref{eq:freshweights95} gives
\begin{equation}\label{eq:scalarcertificate95}
 (T-ezw)^2-D(w)=C(w-w_*)^2.
\end{equation}
For clarity, its quadratic coefficient is
$e^2z^2-(1-d^2v)=C$, its linear coefficient is
$-2z(eT-1)$, and its constant coefficient equals
$z^2(eT-1)^2/C$.  The last identity follows by squaring
$\sqrt{C/(d^2-1)}$ in the expression for $T$.
Since $T\ge e$ and $z<1$, $T-ezw>0$ throughout $[-1,1]$.
Moreover
\[
 eT-1=d(d-2)+
     \frac{e(2d-1)v}{e+d\sqrt{C/(d^2-1)}}\in[0,C],
\]
which proves $0\le w_*\le z<1$.
Taking square roots in \eqref{eq:scalarcertificate95} proves
$F_z(w)\le(e+T)/2=\psi_d(v)$, with equality at $w_*$.
No squared stationary equation with an unchecked sign is used.
The rank-one and zero cases were treated before division by $C$.
Finally polar isometries preserve these trace norms for all rectangular
coefficient maps having Grams $A,B$.
\end{proof}

\begin{lemma}[Concavity of the rank-two spectral objective]
\label{lem:psiconcavity95}
The function $\psi_d$ is continuous, strictly increasing and concave
on $[0,1]$, with
\[
 \psi_d(0)=d-1,\qquad\psi_d(1)=d-1/2.
\]
The function $f_d(x,1-x)=\psi_d(4x(1-x))$ is symmetric and
strictly concave on the two-coordinate probability simplex.
\end{lemma}
\begin{proof}
In the positive added term of \eqref{eq:psidef95} write
\[
 a=d-1,\quad b=d/\sqrt{d^2-1},\quad c=d(d-2),\quad s=2d-1.
\]
It suffices to examine $h(v)=v/(a+b\sqrt{c+sv})$.
For $y=\sqrt{c+sv}>0$ differentiation gives
\[
 h'(v)=\frac{a+(b/2)(y+c/y)}{(a+by)^2}>0.
\]
The derivative of this last expression with respect to $y$ is
\[
 -\frac{b}{2(a+by)^3}
       \left(3a+\frac{ac}{y^2}+by+\frac{3bc}{y}\right)<0.
\]
Since $y$ increases with $v$, this proves concavity in the interior.
When $d=2$ the formula extends continuously to $v=0$ with finite
right derivative; concavity on the closed interval follows by limits.
The endpoint values follow by substitution.  Finally $4x(1-x)$
is strictly concave and symmetric, and composition with an increasing
concave function gives the claimed strict concavity of $f_d$.
\end{proof}

\subsection{Full protocol optimization and attainment}
\begin{proof}[Proof of Theorem~\ref{thm:equalinitial95}]
Write $\alpha=t^2/[2d^2(d+1)]$.
The two equal-prior payoff blocks are those of
\eqref{eq:equalpayoff91}: $\alpha(I-dS)$ for equal labels and
$-\alpha(I-dS)/(d-1)$ for unequal labels.
Use the dimension-preserving normal form of
Lemma~\ref{lem:ranknormal92}, on the prescribed initial acquisition.
For each first label, its Grams satisfy
\[
 \sum_h A_{yh}=\rho,\quad \operatorname{rank}A_{yh}\le2,
 \quad \tr b_{yh}=1,\quad\operatorname{rank}b_{yh}\le\ell.
\]
On a nonzero branch put $w_{yh}=\tr A_{yh}$ and
$a_{yh}=A_{yh}/w_{yh}$.  These are Gram weights, not conditional
history probabilities.  Optimizing final effects over each second
label gives exactly $\alpha$ times the trace norm of the centered
filtered swap: the $d-1$ opposite-sign blocks cancel their denominator,
and $\tr X_++\tr(-X)_+=\|X\|_1$.
Lemma~\ref{lem:ranktwoleaf95} therefore bounds the gain by
\[
 \alpha\sum_{y,h}w_{yh}f_d(\lambda(a_{yh})).
\]
Apply Lemma~\ref{lem:spectralroof94} with rank two and the
continuous symmetric concave function of
Lemma~\ref{lem:psiconcavity95}.  For every $y$ the corresponding
sum is at most $f_d(q,1-q)=\psi_d(v_\rho)$.
There are $d$ first labels.  This proves the upper in
\eqref{eq:equalinitial95}.  Kraus refinements of mixed receiver maps
and classical refinements of fresh mixtures do not retain an extra
quantum environment.  Countable outcomes and null histories use
Lemma~\ref{lem:countable91}; the bound then passes to tester closure.

For attainment choose the at-most-$d$ commuting atoms of
Lemma~\ref{lem:spectralroof94},
$\rho=\sum_hw_h a_h$, all of spectrum $(q,1-q,0,\ldots,0)$.
Factor $w_h a_h=C_h^*C_h$ with $C_h:\C^d\to\C^2$.
The maps $V_h=C_hC_0^+$ satisfy
$V_hC_0=C_h$ and $\sum_hV_h^*V_h=I$ on the initial Schmidt
support, and extend to a complete recorded instrument on the initial
receiver.  For each $h$, choose the commuting maximizing $b_h$
of Lemma~\ref{lem:ranktwoleaf95} and a coefficient map
$D_h:\C^d\to\C^2$ with $D_h^*D_h=b_h$.
Prepare its normalized vector independently of the old output,
conditional on $h$.  For the centered receiver operator $X_h$,
declare $C$ on its positive spectral subspace when labels agree,
and on its negative spectral subspace when they differ; assign the
kernel arbitrarily.  The complementary effects complete a binary
POVM on every block.  Every leaf bound and the envelope bound is now
attained.  The instrument is common to all first labels and is defined
without dividing by any physical branch probability.

The endpoint and strictness claims follow from
Lemma~\ref{lem:psiconcavity95}.  For the rank-one assertion,
Theorem~\ref{thm:rankspectrum92} gives the upper independently of
$\rho$.  Decompose this fixed Gram into pure spectral atoms and
construct the same support-inverse recorded instrument.  Prepare an
aligned pure fresh input for each atom.  Its centered norm is $d-1$,
so the positive-part decisions attain that upper with both receivers
of dimension one.  This proves the assertion whenever either of the
allowed dimensions is one.  The one-call channel identity and
independent-redraw conclusion follow from the first moment, exactly as
in Theorem~\ref{thm:equalprior91}.
\end{proof}

\subsection{A complete qubit formula and spectral inference}
\begin{corollary}[Every pure initial qubit acquisition]
\label{cor:qubitinitial95}
For $d=2$ and any $\rho\in\mathcal D_2$, put
$u=\sqrt{\det\rho}$.  With both retained registers allowed dimension
two, the exact equal-prior reset value is
\begin{equation}\label{eq:qubitinitial95}
 P_{2,2}^{\rm eq}(\rho)
  =\frac12+\frac{t^2}{12}
                     \left(1+\frac{6u^2}{1+4u}\right).
\end{equation}
After identifying the initial Schmidt support with a subspace of $\C^2$,
no nontrivial recorded receiver instrument or receiver message is required.
In an eigenbasis of $\rho$, write its eigenvalues as
$(1+z)/2,(1-z)/2$, $0\le z\le1$.  An optimal fresh Gram has weights
\begin{equation}\label{eq:qubitfresh95}
 \frac12\left(1+\frac{z}{1+2\sqrt{1-z^2}}\right),\qquad
 \frac12\left(1-\frac{z}{1+2\sqrt{1-z^2}}\right).
\end{equation}
The final joint spectral decisions are as in the theorem.
All receiver feedback, public randomization and intermediate receiver
processing are included in the upper bound.
\end{corollary}
\begin{proof}
Here the initial rank is at most two, so its rank-two majorant is its
own spectrum.  In the preceding attainment the single atom $\rho$
suffices: keep the first receiver unchanged on that Schmidt support.  In
\eqref{eq:psidef95}, $v_\rho=4\det\rho=4u^2$ and
$\psi_2(v)=1+3v/[2(1+2\sqrt v)]$.
This gives \eqref{eq:qubitinitial95}.
For $v>0$, \eqref{eq:freshweights95} simplifies to
$w_*=z/(1+2\sqrt v)$; its endpoint limit gives the stated pure
fresh Gram at $z=1$.
\end{proof}
Thus a pure initial Gram gives $1/2+t^2/12$ and the maximally mixed
Gram gives $1/2+t^2/8$.  At a strictly intermediate spectrum the
maximizing fresh spectrum is generally neither the initial spectrum
nor the uniform spectrum.  This differs from the biased rank-two
leaf, whose optimal fresh weights are equal.

\begin{corollary}[Exact loss and a calibrated spectral bound]
\label{cor:equalspectralloss95}
In Theorem~\ref{thm:equalinitial95}, the exact loss from the optimized
rank-two value is
\begin{equation}\label{eq:equalspectralloss95}
 \frac{t^2}{2d(d+1)}
                \left(d-\frac12-\psi_d(v_\rho)\right).
\end{equation}
Suppose $t>0$ and an implemented score $p$ differs by at most
$\eta\ge0$ from the score of a nominal strategy with this fixed
initial acquisition and the stated reset dimensions.  Put
\[
 R=\frac{2d(d+1)}{t^2}\left(p-\eta-\frac12\right).
\]
If $d-1<R\le d-1/2$, let $v_R\in(0,1]$ be the unique solution
$\psi_d(v_R)=R$.  Then
\begin{equation}\label{eq:equalcapinference95}
 \lambda_{\max}(\rho)\le\frac{1+\sqrt{1-v_R}}2.
\end{equation}
If $R\le d-1$ this reasoning gives only the trivial bound one.
If $R>d-1/2$, the nominal model and the supplied error bound cannot
both describe that score.
\end{corollary}
\begin{proof}
Subtract \eqref{eq:equalinitial95} from
\eqref{eq:rankequal92} at rank two.  For the second assertion,
$p-\eta\le P_{2,\ell}^{\rm eq}(\rho)$ gives
$R\le\psi_d(v_\rho)$.  Strict monotonicity gives
$v_\rho\ge v_R$, and solving $4q(1-q)\ge v_R$ for $q\ge1/2$
proves \eqref{eq:equalcapinference95}.  The two remaining cases follow
from the endpoint range of $\psi_d$.
\end{proof}
The error hypothesis concerns a specified implementation and nominal
experiment, with fixed priors and latent-index law.  A uniform
unhalved diamond error $\delta$ on each complete device channel
contributes at most $\delta$ to a two-call decision score; a certified
hard-path preparation error $\varepsilon$ contributes at most
$\varepsilon/2$.  Their sum may therefore be used for $\eta$, as
in Corollary~\ref{cor:gamestability91}.  A fidelity estimate, a sampled
matrix inequality or an uncalibrated score is not substituted for this
hypothesis.  The bound infers a property of the prescribed initial Gram,
not its eigenbasis or the dimension of an unrestricted coherent machine.
